% Section 2 --- Atomicity Taxonomy and System Model (~2.0 pages)
%
% REVIEW STATUS: First complete draft (2026-05-28).
% Evidence sources: Astria docs (TF-3 verbatim), Espresso architecture.md
% (TF-3 verbatim), mev-commit docs (conflation), N=30 survey (§2.4).
% OPEN: §2.3 table --- Radius row pending Task 2.3 codebase review.

\section{Atomicity Taxonomy and System Model}
\label{sec:taxonomy}

\subsection{The Five Layers}
\label{sec:layers}

Sequencing commitments in multi-rollup systems can provide different
levels of ordering and execution guarantee. We distinguish five layers,
A0 through A4, ordered from weakest to strongest.

\begin{table}[h]
\centering
\small
\caption{The A0--A4 atomicity taxonomy. Each layer strictly implies the
  layers below it. ``Sequencing commitment'' refers to any commitment
  issued by a shared sequencer or preconfirmation layer before rollup
  execution.}
\label{tab:taxonomy}
\begin{tabular}{@{}lll@{}}
\toprule
Level & Property & Mechanism closes it \\
\midrule
A0 & \textbf{Common ordering}: all rollups see the same transaction & Any shared sequencer \\
   & order within a bundle commitment. & \\
A1 & \textbf{Atomic inclusion}: both bundle legs are included in the & Joint block / meta block \\
   & same sequencer block or commitment. & \\
A2 & \textbf{Execution atomicity}: all bundle legs succeed, or none & Capability (i) or (ii) \\
   & do. No partial execution persists. & (Theorem~\ref{thm:a2-necessary}) \\
A3 & \textbf{State consistency}: the post-execution states of all & Cross-chain state \\
   & rollups reflect the bundle as a unit. & synchronization \\
A4 & \textbf{Economic atomicity}: no party extracts strictly positive & Bond / slashing / \\
   & option value from the sequencing gap. & insurance (A4 only) \\
\bottomrule
\end{tabular}
\end{table}

\noindent The term ``atomicity'' appears throughout the rollup sequencing
literature without qualification. The Astria shared sequencer documentation
claims \emph{``atomic cross-rollup composability''}~\cite{Astria-README},
and the mev-commit preconfirmation protocol advertises \emph{``Atomicity
and Execution Guarantees''}~\cite{MevCommit-Marketing}. As we show in
\S\ref{sec:obs1} and \S\ref{sec:designspace}, both systems provide at most
A1 and A4 guarantees. The conflation of these with A2 is the structural
gap this paper addresses.

%----------------------------------------------------------------------
\subsection{The A2--A4 Separation}
\label{sec:a2a4}

The central observation of this paper is that A2 and A4 have distinct
closure mechanisms that do not substitute for each other.

\paragraph{A4 closes economic externality; A2 closes execution outcome.}
A4 mechanisms (bonds, slashing, restaking) guarantee that if A2 fails, the
violating party suffers a financial penalty proportional to the violation.
This addresses the \emph{economic consequence} of partial execution but
does not prevent it. An attacker who finds a commitment window in which
$V > B$ --- where $V$ is the option value and $B$ is the bond --- will violate
A2 and accept the slashing.

\paragraph{A2 requires structural capability.}
Theorem~\ref{thm:a2-necessary} (\S\ref{sec:thm1}) establishes that any
deterministic protocol guaranteeing A2 must implement at least one of:
(i) pre-execution success-predicate evidence, or (ii) authority to prevent
partial irreversible effects. Probabilistic and cryptoeconomic guarantees
do not provide either. This is not a quantitative difference; it is a
categorical one.

\paragraph{Preview of Theorem~\ref{thm:bond-insufficient}.}
Theorem~\ref{thm:bond-insufficient} (\S\ref{sec:thm2}) shows that under
a power-law option-value tail, finite bonds produce a strictly positive
expected A4 violation count regardless of bond size. The structural gap
between A4 and A2 therefore has an empirical consequence: even
optimally-sized bonds cannot close it.

%----------------------------------------------------------------------
\subsection{Deployed System Mapping}
\label{sec:mapping}

Table~\ref{tab:mapping} maps the primary shared sequencer and
preconfirmation systems to the A0--A4 taxonomy. All confirmed entries
are derived from primary source documentation read verbatim.

\begin{table}[h]
\centering
\small
\caption{Atomicity level provided by deployed shared sequencer and
  preconfirmation systems. ``Confirmed'' entries cite verbatim primary
  source documentation; ``Conditional'' entries depend on configuration
  or version. $\dagger$See Appendix~\ref{app:boundary}.}
\label{tab:mapping}
\begin{tabular}{@{}lllp{4.5cm}@{}}
\toprule
System & Provides & Does not provide & Key evidence \\
\midrule
Astria shared seq. & A0, A1 & A2, A3, A4 &
  \emph{``provides a guarantee on the ordering of transactions in a
  block, but it doesn't execute the STF of any given rollup''}~\cite{Astria-TxFlow}\\
Espresso HotShot & A0, A1 & A2, A3, A4 &
  \emph{``L2 validators receive blocks and execute the STFs for their
  rollups''} (after HotShot confirmation)~\cite{Espresso-Architecture} \\
mev-commit (Primev) & A0, A1, A4 & A2, A3 &
  \emph{``Array of transaction hashes that can revert''} in bid;
  slashing on positioning failure~\cite{MevCommit-Docs} \\
Bolt (Chainbound) & A0, A1, A4 & A2, A3 &
  Signed preconf on inclusion; no execution outcome
  predicate~\cite{Bolt-Docs} \\
Espresso + Hyperlane$^\dagger$ & A0, A1, A2 (relay) & A3, A4 &
  \texttt{onlyMailbox} modifier; destination executes only after
  Hyperlane relay delivers source-chain
  event~\cite{Espresso-Presto} \\
Radius (encrypted seq.) & A0, A1 & A2, A3, A4 &
  Encrypted transaction ordering; no joint
  execution~\cite{Radius-Docs} \\
\bottomrule
\end{tabular}
\end{table}

\noindent The Espresso + Hyperlane row warrants elaboration. The
\texttt{onlyMailbox} modifier on destination contracts is a cap.~(ii)
implementation: it prevents the destination from executing until the
Hyperlane relay delivers confirmation of the source-chain event. This
is not an atomicity guarantee from the sequencer --- it is an
application-level capability~(ii) built on top of a sequencer that
provides only A0/A1. The sequencer's role is to provide fast
source-chain finality (via Espresso HotShot), which the relay then
uses to trigger destination execution. This is the correct engineering
pattern~(\S\ref{sec:designspace}), and it is precisely what is absent in
the attack scenario we demonstrate~(\S\ref{sec:pefo}).

%----------------------------------------------------------------------
\subsection{The Integration Gap}
\label{sec:gap}

\paragraph{Survey methodology.}
We systematically surveyed $N=30$ cross-rollup applications against the
three structural conditions required for PEFO vulnerability (detailed in
Appendix~\ref{app:survey}). The mandatory filter requires: (TF-1) the
application spans $\geq 2$ rollups under a shared sequencer commitment;
(TF-2) bundle legs execute independently without messaging relay
dependency; and (TF-3) the sequencer is lazy.

\paragraph{Structural finding.}
No deployed application in our survey relies on a shared-sequencer
commitment for cross-rollup execution without adding relay coordination
or escrow protection on the critical cross-rollup path. This is not a
null result. It is a structural finding: TF-2 (independent execution legs
without relay) is the binding constraint, and it is binding because the
ecosystem has independently converged on closing it.

Every cross-rollup application we surveyed that uses a shared sequencer
routes around the A2 gap via one of two mechanisms:

\begin{enumerate}
  \item \textbf{Relay coordination} (capability (ii), application-level):
    Hyperlane, LayerZero, Wormhole, and similar messaging layers enforce
    a sequential source-relay-destination dependency that prevents partial
    execution. Examples: Stargate~\cite{Stargate}, the Espresso+Hyperlane
    reference pattern~\cite{Espresso-Presto}.
  \item \textbf{Solver inventory or oracle escrow} (capability (ii) or
    A4): Solvers absorb partial-fill risk (T3rn, Relay Protocol), or
    origin-chain escrow is held until oracle confirmation (Across Protocol
    via UMA~\cite{Across-Protocol}).
\end{enumerate}

\noindent This convergence is the strongest possible validation of
Theorem~\ref{thm:a2-necessary}'s necessity claim: production engineers
have independently identified that cross-rollup execution requires
capability~(i) or (ii), and have deployed accordingly. The gap
remains live as a forward-looking risk for applications that build
on lazy-sequencer infrastructure without implementing either mechanism ---
precisely the condition we instantiate in \S\ref{sec:pefo}.

%----------------------------------------------------------------------
\subsection{System and Threat Model}
\label{sec:model}

\paragraph{Rollup architecture.}
We consider two rollups $R_1, R_2$ sharing a lazy sequencer $\mathcal{S}$
(Definition~\ref{def:lazy}). Each rollup $R_i$ has local state $\sigma_i$
and independent state transition function $\delta_i$
(Definition~\ref{def:independent}). The sequencer issues a commitment
$\Gamma(B)$ for bundle $B = (\tau_1, \tau_2)$ before either STF executes.
Full formal definitions are in \S\ref{sec:thm1}.

\paragraph{Preconfirmation.}
A preconfirmation is a signed commitment by a validator or proposer to
include a specific transaction in a future block, issued before
execution~\cite{Drake2023-BasedPreconf}. We distinguish three commitment
levels following~\cite{Drake2023-BasedPreconf}: (strongest) post-execution
state root; (medium) execution state diffs; (weakest) transaction inclusion.
The weakest level is A0/A1 only; the strongest approaches capability~(i)
for single-chain L2. We show in \S\ref{sec:thm1} that extending even the
strongest level to multi-rollup bundles requires a \emph{joint} success
predicate covering both STFs.

\paragraph{Determinism scope.}
Theorem~\ref{thm:a2-necessary} applies to \emph{deterministic} protocols.
Probabilistic guarantees (e.g., optimistic execution with fraud proofs) and
cryptoeconomic guarantees (e.g., restaking bonds) are A4 mechanisms.
Non-deterministic execution outcomes do not change the structural analysis:
if the sequencer cannot predict whether both legs will succeed, it cannot
provide capability~(i), and if it lacks state-revert authority, it cannot
provide capability~(ii).

\paragraph{Irreversibility.}
The \emph{irreversibility point} $\pi_i(\tau_i)$ of transaction $\tau_i$
on rollup $R_i$ is the first protocol-internal moment at which the effect
of $\tau_i$ on $\sigma_i$ cannot be undone by any party within the
protocol (Definition~\ref{def:irreversibility}). For optimistic rollups,
this is typically after the fraud-proof window closes. For ZK rollups,
this is after the validity proof is accepted on L1. For the purposes of
the attack, we consider the rollup sequencer's internal commitment to the
execution result as the relevant irreversibility boundary.

\paragraph{Adversary classes.}
We consider three adversary classes:
\begin{itemize}
  \item $\mathcal{A}_{\text{lazy}}$: a sequencer that is lazy with respect
    to one bundle and exploits the commitment for option value.
  \item $\mathcal{A}_{\text{ext}}$: an external party (e.g., MEV
    searcher) who observes the sequencer's commitment and submits a
    conflicting transaction to cause partial execution.
  \item $\mathcal{A}_{\text{solver}}$: a cross-rollup solver who
    submits independent fill transactions, assuming joint execution.
\end{itemize}
In the PEO attack (\S\ref{sec:peo}), $\mathcal{A}_{\text{lazy}}$ is
the relevant class. In the PEFO attack (\S\ref{sec:pefo}),
$\mathcal{A}_{\text{ext}}$ exploits the position of an honest
$\mathcal{A}_{\text{solver}}$.

\paragraph{Out of scope.}
We do not consider: (i) attacks on the consensus protocol of the sequencer
itself; (ii) L1 reorg attacks; (iii) data availability failures; (iv)
validator key compromise. These are orthogonal to the sequencing-commitment
gap we analyze.

\paragraph{Evidence types.}
Following standard security research convention, we use three evidence
types: $\alpha$ (testbed or emulation-based), $\beta$ (historical
data-based counterfactual), and primary source documentation (verbatim
from official protocol documentation, verified via WebFetch or manual
download by the authors).
